% Source: backpropagation_and_chain_rule.tex
\begin{notechapter}
% ---- chapter-local definitions ----
\DeclareDocumentCommand{\subsubsubsection}{sO{#3}m}{
  \IfBooleanTF{#1}{%
    \paragraph*{#3}\mbox{}%
  }{%
    \paragraph[#2]{#3}\mbox{}%
  }%
}
% -----------------------------------
\chapter{Backpropagation and the Chain Rule}\label{ch:backprop-chain-rule}

\section{Chain Rule of Calculus}
The chain rule is a fundamental concept in calculus used to compute the derivative of a composition of functions. In the context of neural networks, it helps us compute gradients for backpropagation. This section introduces the chain rule in both its developed form and matrix notation, showing how the two representations are connected.

\subsection{Chain Rule in \( \mathbb{R}^3 \rightarrow \mathbb{R}^2 \) Example}
Consider a function \( f: \mathbb{R}^3 \rightarrow \mathbb{R}^2 \):
\[
f(x, y, z) = (f_1(x, y, z), f_2(x, y, z)),
\]
where:
\[
f_1(x, y, z) = x^2 + yz, \quad f_2(x, y, z) = \sin(x) + e^{yz}.
\]

Now, suppose we have another function \( g: \mathbb{R}^2 \rightarrow \mathbb{R} \):
\[
g(u, v) = u \cdot v.
\]

We want to compute the derivative of the composition:
\[
h(x, y, z) = g(f_1(x, y, z), f_2(x, y, z)).
\]

\subsubsection{Developed Form of the Chain Rule}
Using the chain rule, the partial derivatives of \( h \) with respect to \( x \), \( y \), and \( z \) are:
\[
\frac{\partial h}{\partial x} = \frac{\partial g}{\partial u} \cdot \frac{\partial f_1}{\partial x} + \frac{\partial g}{\partial v} \cdot \frac{\partial f_2}{\partial x},
\]
\[
\frac{\partial h}{\partial y} = \frac{\partial g}{\partial u} \cdot \frac{\partial f_1}{\partial y} + \frac{\partial g}{\partial v} \cdot \frac{\partial f_2}{\partial y},
\]
\[
\frac{\partial h}{\partial z} = \frac{\partial g}{\partial u} \cdot \frac{\partial f_1}{\partial z} + \frac{\partial g}{\partial v} \cdot \frac{\partial f_2}{\partial z}.
\]

These equations represent the chain rule in its developed form. Each partial derivative of \( h \) is a sum of products, where each product involves a partial derivative of \( g \) and a partial derivative of \( f \).

\subsubsection{Matrix Representation of the Chain Rule}
The chain rule can also be expressed compactly using matrix notation. To do this, we introduce the \textbf{Jacobian matrix}, which is a matrix of all first-order partial derivatives of a vector-valued function.

For \( f: \mathbb{R}^3 \rightarrow \mathbb{R}^2 \), the Jacobian matrix \( J_f \) is:
\[
J_f = \frac{\partial \mathbf{f}}{\partial \mathbf{x}} = \begin{bmatrix}
\frac{\partial f_1}{\partial x} & \frac{\partial f_1}{\partial y} & \frac{\partial f_1}{\partial z} \\
\frac{\partial f_2}{\partial x} & \frac{\partial f_2}{\partial y} & \frac{\partial f_2}{\partial z}
\end{bmatrix}.
\]

\textbf{Understanding the Jacobian Matrix}:
The Jacobian matrix can be thought of as a collection of gradients, where each row corresponds to the gradient of one of the component functions \( f_i \). For example:
\begin{itemize}
    \item The first row \( \begin{bmatrix} \frac{\partial f_1}{\partial x} & \frac{\partial f_1}{\partial y} & \frac{\partial f_1}{\partial z} \end{bmatrix} \) is the gradient of \( f_1 \).
    \item The second row \( \begin{bmatrix} \frac{\partial f_2}{\partial x} & \frac{\partial f_2}{\partial y} & \frac{\partial f_2}{\partial z} \end{bmatrix} \) is the gradient of \( f_2 \).
\end{itemize}

Thus, the Jacobian matrix is essentially a "stack" of gradients, one for each component function \( f_i \).

For the example above, the Jacobian matrix is:
\[
J_f = \begin{bmatrix}
2x & z & y \\
\cos(x) & z e^{yz} & y e^{yz}
\end{bmatrix}.
\]

The gradient of \( h \) with respect to \( \mathbf{x} = [x, y, z]^T \) can be expressed using the Jacobian matrix:
\[
\frac{\partial h}{\partial \mathbf{x}} = \frac{\partial g}{\partial \mathbf{f}} \cdot J_f,
\]
where \( \frac{\partial g}{\partial \mathbf{f}} \) is the \textbf{gradient of \( g \) with respect to \( \mathbf{f} \)}. Since \( g: \mathbb{R}^2 \rightarrow \mathbb{R} \) is a scalar-valued function, \( \frac{\partial g}{\partial \mathbf{f}} \) is a row vector:
\[
\frac{\partial g}{\partial \mathbf{f}} = \begin{bmatrix} \frac{\partial g}{\partial f_1} & \frac{\partial g}{\partial f_2} \end{bmatrix}.
\]

For \( g(u, v) = u \cdot v \), we have:
\[
\frac{\partial g}{\partial f_1} = f_2 \quad \text{and} \quad \frac{\partial g}{\partial f_2} = f_1.
\]
Thus:
\[
\frac{\partial g}{\partial \mathbf{f}} = \begin{bmatrix} f_2 & f_1 \end{bmatrix}.
\]

The gradient of \( h \) is then:
\[
\frac{\partial h}{\partial \mathbf{x}} = \begin{bmatrix} f_2 & f_1 \end{bmatrix} \cdot \begin{bmatrix}
2x & z & y \\
\cos(x) & z e^{yz} & y e^{yz}
\end{bmatrix}.
\]

\textbf{Connection to the Developed Form}:
The matrix multiplication \( \frac{\partial g}{\partial \mathbf{f}} \cdot J_f \) is equivalent to the chain rule applied in the developed form. For example, the first entry of the resulting gradient vector is:
\[
\frac{\partial h}{\partial x} = f_2 \cdot \frac{\partial f_1}{\partial x} + f_1 \cdot \frac{\partial f_2}{\partial x},
\]
which matches the expression for \( \frac{\partial h}{\partial x} \) derived using the chain rule. Similarly, the other entries of the gradient vector correspond to the partial derivatives with respect to \( y \) and \( z \).

\subsection{Gradient Convention}
The gradient of a scalar-valued function can be defined as either a row vector or a column vector, depending on the convention:
\begin{itemize}
    \item \textbf{Row Vector Convention}:
    \[
    \nabla f = \begin{bmatrix} \frac{\partial f}{\partial x_1} & \frac{\partial f}{\partial x_2} & \cdots & \frac{\partial f}{\partial x_n} \end{bmatrix}.
    \]
    This is commonly used in optimization and machine learning.

    \item \textbf{Column Vector Convention}:
    \[
    \nabla f = \begin{bmatrix} \frac{\partial f}{\partial x_1} \\ \frac{\partial f}{\partial x_2} \\ \vdots \\ \frac{\partial f}{\partial x_n} \end{bmatrix}.
    \]
    This is often used in physics and engineering.
\end{itemize}

In this document, we use the \textbf{row vector convention} for consistency with the chain rule as applied in backpropagation.

\section{Backpropagation in Neural Networks}
Backpropagation is the process of computing gradients in a neural network using the chain rule. It consists of a forward pass and a backward pass.

\subsection{Forward Pass}
In the forward pass, we compute the output of the network by applying layer-by-layer transformations. For a 2-layer fully connected MLP:
\begin{itemize}
    \item Input: \( \mathbf{x} \in \mathbb{R}^d \),
    \item Hidden layer: \( \mathbf{h} = \sigma(W_1 \mathbf{x} + \mathbf{b}_1) \),
    \item Output: \( \mathbf{y} = W_2 \mathbf{h} + \mathbf{b}_2 \),
    \item Loss: \( \mathcal{L} \).
\end{itemize}

Here, \( W_1 \in \mathbb{R}^{m \times d} \), \( W_2 \in \mathbb{R}^{n \times m} \), \( \mathbf{b}_1 \in \mathbb{R}^m \), \( \mathbf{b}_2 \in \mathbb{R}^n \), and \( \sigma \) is an element-wise activation function.

\subsection{Backward Pass}
In the backward pass, we compute gradients layer-by-layer, starting from the output and moving backward to the input.

\subsubsection{Gradients for Output Layer}
The gradient of the loss \( \mathcal{L} \) with respect to \( W_2 \) is:
\[
\frac{\partial \mathcal{L}}{\partial W_2} = \left( \frac{\partial \mathcal{L}}{\partial \mathbf{y}} \right)^T \cdot \mathbf{h}^T.
\]

Here:
\begin{itemize}
    \item \( \frac{\partial \mathcal{L}}{\partial \mathbf{y}} \) is a row vector of size \( 1 \times n \),
    \item \( \mathbf{h}^T \) is a row vector of size \( 1 \times m \),
    \item The product \( \left( \frac{\partial \mathcal{L}}{\partial \mathbf{y}} \right)^T \cdot \mathbf{h}^T \) is a matrix of size \( n \times m \), which matches the dimensions of \( W_2 \).
\end{itemize}

\textbf{Explanation}:
The gradient \( \frac{\partial \mathcal{L}}{\partial W_2} \) is computed using the outer product of \( \frac{\partial \mathcal{L}}{\partial \mathbf{y}} \) and \( \mathbf{h}^T \). This ensures that the dimensions are consistent and the gradient is computed correctly.

\subsubsection{Detailed Example with \( W_2 \) as a \( 2 \times 3 \) Matrix}
To better understand why \( \mathbf{h}^T \) appears in the gradient computation, let's work through an example where \( W_2 \) is a \( 2 \times 3 \) matrix.

\subsubsection{Setup}
Let \( W_2 \) be a \( 2 \times 3 \) matrix:
\[
W_2 = \begin{bmatrix}
w_{11} & w_{12} & w_{13} \\
w_{21} & w_{22} & w_{23}
\end{bmatrix}.
\]
Let \( \mathbf{h} \) be a \( 3 \times 1 \) vector:
\[
\mathbf{h} = \begin{bmatrix}
h_1 \\
h_2 \\
h_3
\end{bmatrix}.
\]
The output \( \mathbf{y} \) is computed as:
\[
\mathbf{y} = W_2 \mathbf{h} + \mathbf{b}_2,
\]
where \( \mathbf{b}_2 \) is a \( 2 \times 1 \) bias vector:
\[
\mathbf{b}_2 = \begin{bmatrix}
b_1 \\
b_2
\end{bmatrix}.
\]
Thus, \( \mathbf{y} \) is a \( 2 \times 1 \) vector:
\[
\mathbf{y} = \begin{bmatrix}
y_1 \\
y_2
\end{bmatrix}.
\]

\subsubsection{Expand \( \mathbf{y} \) in Terms of \( W_2 \)}
The output \( \mathbf{y} \) is given by:
\[
\mathbf{y} = W_2 \mathbf{h} + \mathbf{b}_2.
\]
Expanding this, we get:
\[
y_1 = w_{11} h_1 + w_{12} h_2 + w_{13} h_3 + b_1,
\]
\[
y_2 = w_{21} h_1 + w_{22} h_2 + w_{23} h_3 + b_2.
\]

\subsubsection{Compute \( \frac{\partial \mathbf{y}}{\partial W_2} \)}
The derivative \( \frac{\partial \mathbf{y}}{\partial W_2} \) is a 3D tensor because \( \mathbf{y} \) is a \( 2 \times 1 \) vector and \( W_2 \) is a \( 2 \times 3 \) matrix. Specifically, \( \frac{\partial \mathbf{y}}{\partial W_2} \) has dimensions \( 2 \times 2 \times 3 \), where:
\begin{itemize}
    \item The first dimension corresponds to the elements of \( \mathbf{y} \) (\( y_1 \) and \( y_2 \)),
    \item The second and third dimensions correspond to the rows and columns of \( W_2 \) (\( 2 \times 3 \)).
\end{itemize}

\subparagraph*{Step 1: Compute \( \frac{\partial y_1}{\partial W_2} \)}
The derivative of \( y_1 \) with respect to \( W_2 \) is a \( 2 \times 3 \) matrix:
\[
\frac{\partial y_1}{\partial W_2} = \begin{bmatrix}
\frac{\partial y_1}{\partial w_{11}} & \frac{\partial y_1}{\partial w_{12}} & \frac{\partial y_1}{\partial w_{13}} \\
\frac{\partial y_1}{\partial w_{21}} & \frac{\partial y_1}{\partial w_{22}} & \frac{\partial y_1}{\partial w_{23}}
\end{bmatrix}.
\]

Since \( y_1 \) depends only on the first row of \( W_2 \), we have:
\[
\frac{\partial y_1}{\partial w_{11}} = h_1, \quad \frac{\partial y_1}{\partial w_{12}} = h_2, \quad \frac{\partial y_1}{\partial w_{13}} = h_3,
\]
and all other derivatives are zero. Thus:
\[
\frac{\partial y_1}{\partial W_2} = \begin{bmatrix}
h_1 & h_2 & h_3 \\
0 & 0 & 0
\end{bmatrix}.
\]

\subparagraph*{Step 2: Compute \( \frac{\partial y_2}{\partial W_2} \)}
Similarly, the derivative of \( y_2 \) with respect to \( W_2 \) is:
\[
\frac{\partial y_2}{\partial W_2} = \begin{bmatrix}
0 & 0 & 0 \\
h_1 & h_2 & h_3
\end{bmatrix}.
\]

\subparagraph*{Step 3: Combine the Results}
The full derivative \( \frac{\partial \mathbf{y}}{\partial W_2} \) is a 3D tensor with slices corresponding to \( \frac{\partial y_1}{\partial W_2} \) and \( \frac{\partial y_2}{\partial W_2} \).

\subsubsection{Compute \( \frac{\partial \mathcal{L}}{\partial W_2} \)}
Using the chain rule, the gradient of the loss \( \mathcal{L} \) with
respect to \( W_2 \) is:
\[
\frac{\partial \mathcal{L}}{\partial W_2} = \left( \frac{\partial \mathcal{L}}{\partial \mathbf{y}} \right)^T \cdot \mathbf{h}^T.
\]
This matches the earlier expression and shows how \( \mathbf{h}^T \) naturally arises in the gradient computation.

\textbf{Treating \( W_2 \) as a Vector}:
When computing gradients, matrices like \( W_2 \) are often treated as flattened vectors. For example, if:
\[
W_2 = \begin{bmatrix}
w_{11} & w_{12} & \cdots & w_{1m} \\
w_{21} & w_{22} & \cdots & w_{2m} \\
\vdots & \vdots & \ddots & \vdots \\
w_{n1} & w_{n2} & \cdots & w_{nm}
\end{bmatrix},
\]
then the flattened version of \( W_2 \) is:
\[
\text{vec}(W_2) = \begin{bmatrix}
w_{11} \\ w_{12} \\ \vdots \\ w_{1m} \\ w_{21} \\ w_{22} \\ \vdots \\ w_{nm}
\end{bmatrix}.
\]
This allows us to treat \( W_2 \) as a vector of length \( n \times m \), and the gradient \( \frac{\partial \mathcal{L}}{\partial W_2} \) is also a matrix of the same size as \( W_2 \), where each entry is the partial derivative of \( \mathcal{L} \) with respect to the corresponding element of \( W_2 \).

\textbf{Why Treat \( W_2 \) as a Vector?}
Treating \( W_2 \) as a vector is useful because:
\begin{itemize}
    \item \textbf{Gradient Descent}: In optimization algorithms like gradient descent, we update the parameters (e.g., \( W_2 \)) by subtracting a scaled version of the gradient. This requires the gradient to have the same shape as the parameter, so treating \( W_2 \) as a vector simplifies the update rule.
    \item \textbf{Chain Rule}: When applying the chain rule, it is often easier to work with vectors and matrices rather than tensors (higher-dimensional arrays). By flattening \( W_2 \), we can use standard vector calculus rules.
    \item \textbf{Implementation}: Many deep learning frameworks (e.g., PyTorch, TensorFlow) internally treat matrices as vectors for gradient computations, making the implementation more efficient.
\end{itemize}

\subsubsection{Gradients for Hidden Layer}
The gradient of the loss \( \mathcal{L} \) with respect to the hidden layer activations \( \mathbf{h} \) is a row vector:
\[
\frac{\partial \mathcal{L}}{\partial \mathbf{h}} = \begin{bmatrix} 
\frac{\partial \mathcal{L}}{\partial h_1} & \frac{\partial \mathcal{L}}{\partial h_2} 
& \cdots & \frac{\partial \mathcal{L}}{\partial h_m} \end{bmatrix}.
\]

This row vector represents the sensitivity of the loss \( \mathcal{L} \) to changes in each element of \( \mathbf{h} \).

\subsubsection{Gradient with respect to \( W_1 \)}
The gradient of the loss \( \mathcal{L} \) with respect to \( W_1 \) is:
\[
\frac{\partial \mathcal{L}}{\partial W_1} = \left( \frac{\partial \mathcal{L}}{\partial \mathbf{h}} 
\odot \sigma'(\mathbf{z}) \right)^T \cdot \mathbf{x}^T,
\]
where:
\begin{itemize}
    \item \( \frac{\partial \mathcal{L}}{\partial \mathbf{h}} \) is a row vector of size \( 1 \times m \),
    \item \( \sigma'(\mathbf{z}) \) is a column vector of size \( m \times 1 \),
    \item \( \mathbf{x}^T \) is a row vector of size \( 1 \times d \).
\end{itemize}

\subsubsection{Gradients for Hidden Layer}
The gradient of the loss \( \mathcal{L} \) with respect to the hidden layer parameters \( W_1 \) and \( \mathbf{b}_1 \) is computed similarly. For example:
\[
\frac{\partial \mathcal{L}}{\partial W_1} = \frac{\partial \mathcal{L}}{\partial \mathbf{h}} \odot \sigma'(W_1 \mathbf{x} + \mathbf{b}_1) \cdot \mathbf{x}^T.
\]
Here, \( \odot \) denotes the \textbf{Hadamard product} (element-wise multiplication), which arises because the activation function \( \sigma \) is applied element-wise.

\subsection{Detailed Development of Gradients}

\subsubsection{Gradient with respect to \( W_2 \)}

The output \( \mathbf{y} \) is computed as:
\[
\mathbf{y} = W_2 \mathbf{h} + \mathbf{b}_2,
\]
where:
\begin{itemize}
    \item \( W_2 \) is an \( n \times m \) weight matrix,
    \item \( \mathbf{h} \) is an \( m \times 1 \) hidden layer activation vector,
    \item \( \mathbf{b}_2 \) is an \( n \times 1 \) bias vector,
    \item \( \mathbf{y} \) is an \( n \times 1 \) output vector.
\end{itemize}

The gradient of the loss \( \mathcal{L} \) with respect to \( W_2 \) is computed using the chain rule:
\[
\frac{\partial \mathcal{L}}{\partial W_2} = \frac{\partial \mathcal{L}}{\partial \mathbf{y}} \cdot \frac{\partial \mathbf{y}}{\partial W_2}.
\]

Here:
\begin{itemize}
    \item \( \frac{\partial \mathcal{L}}{\partial \mathbf{y}} \) is a \( 1 \times n \) row vector (the gradient of the scalar loss \( \mathcal{L} \) with respect to the vector \( \mathbf{y} \)),
    \item \( \frac{\partial \mathbf{y}}{\partial W_2} \) is a 3D tensor of size \( n \times n \times m \) (the derivative of the vector \( \mathbf{y} \) with respect to the matrix \( W_2 \)).
\end{itemize}

However, working with 3D tensors is cumbersome, so we instead compute the gradient element-wise and assemble the result into a matrix.

\subparagraph*{Step 1: Element-Wise Gradient Computation}
Let’s compute the gradient of \( \mathcal{L} \) with respect to a single element \( w_{ij} \) of \( W_2 \), where \( w_{ij} \) is the element in the \( i \)-th row and \( j \)-th column of \( W_2 \).

From \( \mathbf{y} = W_2 \mathbf{h} + \mathbf{b}_2 \), the \( i \)-th element of \( \mathbf{y} \) is:
\[
y_i = \sum_{k=1}^m w_{ik} h_k + b_i.
\]

The derivative of \( y_i \) with respect to \( w_{ij} \) is:
\[
\frac{\partial y_i}{\partial w_{ij}} = h_j.
\]

Thus, the derivative of \( \mathcal{L} \) with respect to \( w_{ij} \) is:
\[
\frac{\partial \mathcal{L}}{\partial w_{ij}} = \frac{\partial \mathcal{L}}{\partial y_i} \cdot \frac{\partial y_i}{\partial w_{ij}} = \frac{\partial \mathcal{L}}{\partial y_i} \cdot h_j.
\]

\subparagraph*{Step 2: Assemble the Gradient Matrix}

The gradient \( \frac{\partial \mathcal{L}}{\partial W_2} \) is a matrix of the same size as \( W_2 \) (i.e., \( n \times m \)), where each element is given by:
\[
\left( \frac{\partial \mathcal{L}}{\partial W_2} \right)_{ij} = \frac{\partial \mathcal{L}}{\partial y_i} \cdot h_j.
\]

This can be written in matrix form as:
\[
\frac{\partial \mathcal{L}}{\partial W_2} = \left( \frac{\partial \mathcal{L}}{\partial \mathbf{y}} \right)^T \cdot \mathbf{h}^T,
\]
where:
\begin{itemize}
    \item \( \frac{\partial \mathcal{L}}{\partial \mathbf{y}} \) is a \( 1 \times n \) row vector,
    \item \( \left( \frac{\partial \mathcal{L}}{\partial \mathbf{y}} \right)^T \) is an \( n \times 1 \) column vector,
    \item \( \mathbf{h}^T \) is a \( 1 \times m \) row vector,
    \item The product \( \left( \frac{\partial \mathcal{L}}{\partial \mathbf{y}} \right)^T \cdot \mathbf{h}^T \) is an \( n \times m \) matrix.
\end{itemize}

\subparagraph*{Step 3: Why This Works}

The key insight is that the gradient \( \frac{\partial \mathcal{L}}{\partial W_2} \) is the \textbf{outer product} of \( \left( \frac{\partial \mathcal{L}}{\partial \mathbf{y}} \right)^T \) and \( \mathbf{h}^T \). This is because:
\begin{itemize}
    \item Each element \( \left( \frac{\partial \mathcal{L}}{\partial W_2} \right)_{ij} \) depends only on \( \frac{\partial \mathcal{L}}{\partial y_i} \) (the \( i \)-th component of \( \frac{\partial \mathcal{L}}{\partial \mathbf{y}} \)) and \( h_j \) (the \( j \)-th component of \( \mathbf{h} \)),
    \item The outer product \( \left( \frac{\partial \mathcal{L}}{\partial \mathbf{y}} \right)^T \cdot \mathbf{h}^T \) captures all such pairwise products.
\end{itemize}

\subparagraph*{Step 4: Final Expression}

Thus, the rigorous and mathematically sound expression for the gradient of \( \mathcal{L} \) with respect to \( W_2 \) is:
\[
\frac{\partial \mathcal{L}}{\partial W_2} = \left( \frac{\partial \mathcal{L}}{\partial \mathbf{y}} \right)^T \cdot \mathbf{h}^T.
\]

\subparagraph*{Summary of Steps}

1. We start with the chain rule:
   \[
   \frac{\partial \mathcal{L}}{\partial W_2} = \frac{\partial \mathcal{L}}{\partial \mathbf{y}} \cdot \frac{\partial \mathbf{y}}{\partial W_2}.
   \]

2. We compute the derivative element-wise:
   \[
   \frac{\partial \mathcal{L}}{\partial w_{ij}} = \frac{\partial \mathcal{L}}{\partial y_i} \cdot h_j.
   \]

3. We assemble the gradient matrix using the outer product:
   \[
   \frac{\partial \mathcal{L}}{\partial W_2} = \left( \frac{\partial \mathcal{L}}{\partial \mathbf{y}} \right)^T \cdot \mathbf{h}^T.
   \]

This derivation avoids treating matrices as vectors and instead uses proper matrix calculus and the chain rule. It is not a "hack" but a rigorous application of matrix derivatives.
\subsubsection{Gradient with respect to \( \mathbf{h} \)}
The hidden layer activations \( \mathbf{h} \) are computed as:
\[
\mathbf{h} = \sigma(W_1 \mathbf{x} + \mathbf{b}_1).
\]
The gradient of the loss \( \mathcal{L} \) with respect to \( \mathbf{h} \) is:
\[
\frac{\partial \mathcal{L}}{\partial \mathbf{h}} = \frac{\partial \mathcal{L}}{\partial \mathbf{y}} \cdot W_2,
\]

where:
\begin{itemize}
    \item \( \frac{\partial \mathcal{L}}{\partial \mathbf{y}} \) is a row vector of size \( 1 \times n \),
    \item \( W_2 \) is a matrix of size \( n \times m \),
    \item The product \( \frac{\partial \mathcal{L}}{\partial \mathbf{y}} \cdot W_2 \) results in a row vector of size \( 1 \times m \), which matches the dimensions of \( \mathbf{h} \).
\end{itemize}

\textbf{Explanation}:
The gradient \( \frac{\partial \mathcal{L}}{\partial \mathbf{h}} \) is computed using the chain rule, where \( \frac{\partial \mathbf{y}}{\partial \mathbf{h}} = W_2 \) is the Jacobian matrix of the linear transformation \( \mathbf{y} = W_2 \mathbf{h} + \mathbf{b}_2 \). Thus, the gradient of \( \mathcal{L} \) with respect to \( \mathbf{h} \) is the product of the gradient of \( \mathcal{L} \) with respect to \( \mathbf{y} \) and the Jacobian matrix \( W_2 \).

\subsubsection{Explicit Development of \( \frac{\partial \mathcal{L}}{\partial \mathbf{y}} \cdot W_2 \)}
The gradient of the loss \( \mathcal{L} \) with respect to \( \mathbf{h} \) is computed as:
\[
\frac{\partial \mathcal{L}}{\partial \mathbf{h}} = \frac{\partial \mathcal{L}}{\partial \mathbf{y}} \cdot W_2,
\]
where:
\begin{itemize}
    \item \( \frac{\partial \mathcal{L}}{\partial \mathbf{y}} = \begin{bmatrix} 
    \frac{\partial \mathcal{L}}{\partial y_1} & \frac{\partial \mathcal{L}}{\partial y_2} 
    & \cdots & \frac{\partial \mathcal{L}}{\partial y_n} \end{bmatrix} \) 
    (a row vector of size \( 1 \times n \)),
    \item \( W_2 = \begin{bmatrix} 
    w_{11} & w_{12} & \cdots & w_{1m} \\
    w_{21} & w_{22} & \cdots & w_{2m} \\
    \vdots & \vdots & \ddots & \vdots \\
    w_{n1} & w_{n2} & \cdots & w_{nm}
    \end{bmatrix} \) (a matrix of size \( n \times m \)),
    \item The product \( \frac{\partial \mathcal{L}}{\partial \mathbf{y}} \cdot W_2 \) 
    results in a row vector of size \( 1 \times m \), which matches the dimensions of \( \mathbf{h} \).
\end{itemize}

\textbf{Explicit Matrix Multiplication}:
\[
\frac{\partial \mathcal{L}}{\partial \mathbf{h}} = \begin{bmatrix} 
\frac{\partial \mathcal{L}}{\partial y_1} & \frac{\partial \mathcal{L}}{\partial y_2} 
& \cdots & \frac{\partial \mathcal{L}}{\partial y_n} 
\end{bmatrix} \cdot \begin{bmatrix}
w_{11} & w_{12} & \cdots & w_{1m} \\
w_{21} & w_{22} & \cdots & w_{2m} \\
\vdots & \vdots & \ddots & \vdots \\
w_{n1} & w_{n2} & \cdots & w_{nm}
\end{bmatrix}
\]
\[
= \begin{bmatrix}
\sum_{i=1}^n \frac{\partial \mathcal{L}}{\partial y_i} w_{i1} & 
\sum_{i=1}^n \frac{\partial \mathcal{L}}{\partial y_i} w_{i2} & 
\cdots & 
\sum_{i=1}^n \frac{\partial \mathcal{L}}{\partial y_i} w_{im}
\end{bmatrix}.
\]

This row vector represents the gradient of the loss \( \mathcal{L} \) with respect to each element of \( \mathbf{h} \).

\subsubsection{Gradient with respect to \( W_1 \)}
The gradient of the loss \( \mathcal{L} \) with respect to \( W_1 \) is:
\[
\frac{\partial \mathcal{L}}{\partial W_1} = \frac{\partial \mathcal{L}}{\partial \mathbf{h}} \odot \sigma'(W_1 \mathbf{x} + \mathbf{b}_1) \cdot \mathbf{x}^T.
\]

\textbf{Where Does the Hadamard Product Come From?}

The key here is the chain rule. The hidden layer activations \( \mathbf{h} \) are computed as:
\[
\mathbf{h} = \sigma(W_1 \mathbf{x} + \mathbf{b}_1).
\]
To compute \( \frac{\partial \mathcal{L}}{\partial W_1} \), we need to apply the chain rule through the activation function \( \sigma \). The derivative of \( \mathbf{h} \) with respect to \( W_1 \) involves the derivative of \( \sigma \), which is applied element-wise. This is why the Hadamard product \( \odot \) appears.

\textbf{Explicit Development}:
\begin{itemize}
    \item \( \frac{\partial \mathcal{L}}{\partial \mathbf{h}} = \begin{bmatrix} 
    \frac{\partial \mathcal{L}}{\partial h_1} \\ \frac{\partial \mathcal{L}}{\partial h_2} \\ \vdots \\ \frac{\partial \mathcal{L}}{\partial h_m} \end{bmatrix} \) 
    (a column vector of size \( m \times 1 \))
    \item \( \sigma'(W_1 \mathbf{x} + \mathbf{b}_1) = \begin{bmatrix} \sigma'(z_1) \\ \sigma'(z_2) \\ \vdots \\ \sigma'(z_m) \end{bmatrix} \) (a column vector of size \( m \times 1 \)), where \( z_i = \sum_{j=1}^d w_{ij} x_j + b_{i1} \)
    \item The Hadamard product \( \frac{\partial \mathcal{L}}{\partial \mathbf{h}} \odot \sigma'(W_1 \mathbf{x} + \mathbf{b}_1) \) is a column vector of size \( m \times 1 \):
    \[
    \begin{bmatrix} \frac{\partial \mathcal{L}}{\partial h_1} \sigma'(z_1) \\ \frac{\partial \mathcal{L}}{\partial h_2} \sigma'(z_2) \\ \vdots \\ \frac{\partial \mathcal{L}}{\partial h_m} \sigma'(z_m) \end{bmatrix}
    \]
    \item \( \mathbf{x}^T = \begin{bmatrix} x_1 & x_2 & \cdots & x_d \end{bmatrix} \) 
    (a row vector of size \( 1 \times d \))
    \item The product \( \left( \frac{\partial \mathcal{L}}{\partial \mathbf{h}} \odot \sigma'(W_1 \mathbf{x} + \mathbf{b}_1) \right) \cdot \mathbf{x}^T \) results in a matrix of size \( m \times d \), which matches the dimensions of \( W_1 \).
\end{itemize}

\textbf{Explicit Matrix Multiplication}:
\[
\frac{\partial \mathcal{L}}{\partial W_1} = \begin{bmatrix} \frac{\partial \mathcal{L}}{\partial h_1} \sigma'(z_1) \\ \frac{\partial \mathcal{L}}{\partial h_2} \sigma'(z_2) \\ \vdots \\ \frac{\partial \mathcal{L}}{\partial h_m} \sigma'(z_m) \end{bmatrix} \cdot \begin{bmatrix} x_1 & x_2 & \cdots & x_d \end{bmatrix}
\]
\[
= \begin{bmatrix}
\frac{\partial \mathcal{L}}{\partial h_1} \sigma'(z_1) x_1 & \frac{\partial \mathcal{L}}{\partial h_1} \sigma'(z_1) x_2 & \cdots & \frac{\partial \mathcal{L}}{\partial h_1} \sigma'(z_1) x_d \\
\frac{\partial \mathcal{L}}{\partial h_2} \sigma'(z_2) x_1 & \frac{\partial \mathcal{L}}{\partial h_2} \sigma'(z_2) x_2 & \cdots & \frac{\partial \mathcal{L}}{\partial h_2} \sigma'(z_2) x_d \\
\vdots & \vdots & \ddots & \vdots \\
\frac{\partial \mathcal{L}}{\partial h_m} \sigma'(z_m) x_1 & \frac{\partial \mathcal{L}}{\partial h_m} \sigma'(z_m) x_2 & \cdots & \frac{\partial \mathcal{L}}{\partial h_m} \sigma'(z_m) x_d
\end{bmatrix}.
\]

This matrix represents the gradient of the loss \( \mathcal{L} \) with respect to each element of \( W_1 \).

\subsubsection{Detailed Example with \( W_1 \) as a \( 2 \times 3 \) Matrix}
To better understand the gradient computation for \( W_1 \), let's work through an example where \( W_1 \) is a \( 2 \times 3 \) matrix.

\subsubsubsection*{1. Setup}
Let \( W_1 \) be a \( 2 \times 3 \) matrix:
\[
W_1 = \begin{bmatrix}
w_{11} & w_{12} & w_{13} \\
w_{21} & w_{22} & w_{23}
\end{bmatrix}.
\]

Let \( \mathbf{x} \) be a \( 3 \times 1 \) input vector:
\[
\mathbf{x} = \begin{bmatrix}
x_1 \\
x_2 \\
x_3
\end{bmatrix}.
\]

The hidden layer activations \( \mathbf{h} \) are computed as:
\[
\mathbf{h} = \sigma(W_1 \mathbf{x} + \mathbf{b}_1),
\]

where \( \mathbf{b}_1 \) is a \( 2 \times 1 \) bias vector:
\[
\mathbf{b}_1 = \begin{bmatrix}
b_1 \\
b_2
\end{bmatrix}.
\]

Thus, \( \mathbf{h} \) is a \( 2 \times 1 \) vector:
\[
\mathbf{h} = \begin{bmatrix}
h_1 \\
h_2
\end{bmatrix}.
\]

\subsubsubsection*{2. Expand \( \mathbf{h} \) in Terms of \( W_1 \)}
The hidden layer activations \( \mathbf{h} \) are given by:
\[
\mathbf{h} = \sigma(W_1 \mathbf{x} + \mathbf{b}_1).
\]

Expanding this, we get:
\[
h_1 = \sigma(w_{11} x_1 + w_{12} x_2 + w_{13} x_3 + b_1),
\]
\[
h_2 = \sigma(w_{21} x_1 + w_{22} x_2 + w_{23} x_3 + b_2).
\]

\subsubsubsection*{3. Compute \( \frac{\partial \mathbf{h}}{\partial W_1} \)}
The derivative \( \frac{\partial \mathbf{h}}{\partial W_1} \) is a 3D tensor because \( \mathbf{h} \) is a \( 2 \times 1 \) vector and \( W_1 \) is a \( 2 \times 3 \) matrix. Specifically, \( \frac{\partial \mathbf{h}}{\partial W_1} \) has dimensions \( 2 \times 2 \times 3 \), where:
\begin{itemize}
    \item The first dimension corresponds to the elements of \( \mathbf{h} \) (\( h_1 \) and \( h_2 \)),
    \item The second and third dimensions correspond to the rows and columns of \( W_1 \) (\( 2 \times 3 \)).
\end{itemize}

\subparagraph*{Step 1: Compute \( \frac{\partial h_1}{\partial W_1} \)}

The derivative of \( h_1 \) with respect to \( W_1 \) is a \( 2 \times 3 \) matrix:
\[
\frac{\partial h_1}{\partial W_1} = \begin{bmatrix}
\frac{\partial h_1}{\partial w_{11}} & \frac{\partial h_1}{\partial w_{12}} & \frac{\partial h_1}{\partial w_{13}} \\
\frac{\partial h_1}{\partial w_{21}} & \frac{\partial h_1}{\partial w_{22}} & \frac{\partial h_1}{\partial w_{23}}
\end{bmatrix}.
\]

Since \( h_1 \) depends only on the first row of \( W_1 \), we have:
\[
\frac{\partial h_1}{\partial w_{11}} = \sigma'(z_1) x_1, \quad \frac{\partial h_1}{\partial w_{12}} = \sigma'(z_1) x_2, \quad \frac{\partial h_1}{\partial w_{13}} = \sigma'(z_1) x_3,
\]
and all other derivatives are zero. Thus:
\[
\frac{\partial h_1}{\partial W_1} = \begin{bmatrix}
\sigma'(z_1) x_1 & \sigma'(z_1) x_2 & \sigma'(z_1) x_3 \\
0 & 0 & 0
\end{bmatrix}.
\]

\subparagraph*{Step 2: Compute \( \frac{\partial h_2}{\partial W_1} \)}

Similarly, the derivative of \( h_2 \) with respect to \( W_1 \) is:
\[
\frac{\partial h_2}{\partial W_1} = \begin{bmatrix}
0 & 0 & 0 \\
\sigma'(z_2) x_1 & \sigma'(z_2) x_2 & \sigma'(z_2) x_3
\end{bmatrix}.
\]

\subparagraph*{Step 3: Combine the Results}
The full derivative \( \frac{\partial \mathbf{h}}{\partial W_1} \) is a 3D tensor with slices corresponding to \( \frac{\partial h_1}{\partial W_1} \) and \( \frac{\partial h_2}{\partial W_1} \).

\subsubsection{Compute \( \frac{\partial \mathcal{L}}{\partial W_1} \)}
Using the chain rule, the gradient of the loss \( \mathcal{L} \) with respect to \( W_1 \) is:
\[
\frac{\partial \mathcal{L}}{\partial W_1} = \frac{\partial \mathcal{L}}{\partial \mathbf{h}} \odot \sigma'(W_1 \mathbf{x} + \mathbf{b}_1) \cdot \mathbf{x}^T.
\]
This matches the earlier expression and shows how the Hadamard product and the chain rule are applied in the gradient computation.

\subsection{Detailed Explanation of the Hadamard Product in \( \frac{\partial \mathcal{L}}{\partial W_1} \)}

The Hadamard product arises because of the \textbf{element-wise application of the activation function} \( \sigma \) in the hidden layer. When computing the gradient of the loss \( \mathcal{L} \) with respect to \( W_1 \), we need to account for the fact that the activation function \( \sigma \) is applied element-wise to the pre-activation values \( \mathbf{z} = W_1 \mathbf{x} + \mathbf{b}_1 \). This means that the derivative of \( \mathbf{h} \) with respect to \( W_1 \) involves the derivative of \( \sigma \), which is applied element-wise.

Let’s break this down step by step.

\subsubsection{Forward Pass}
The hidden layer activations \( \mathbf{h} \) are computed as:
\[
\mathbf{h} = \sigma(W_1 \mathbf{x} + \mathbf{b}_1),
\]
where:
\begin{itemize}
    \item \( W_1 \) is a \( 2 \times 3 \) matrix,
    \item \( \mathbf{x} \) is a \( 3 \times 1 \) input vector,
    \item \( \mathbf{b}_1 \) is a \( 2 \times 1 \) bias vector,
    \item \( \sigma \) is an element-wise activation function (e.g., ReLU, sigmoid).
\end{itemize}

Let \( \mathbf{z} = W_1 \mathbf{x} + \mathbf{b}_1 \). Then:
\[
\mathbf{h} = \sigma(\mathbf{z}),
\]
where \( \mathbf{z} = \begin{bmatrix} z_1 \\ z_2 \end{bmatrix} \) and \( \mathbf{h} = \begin{bmatrix} h_1 \\ h_2 \end{bmatrix} \).

\subsubsection{Backward Pass}
The gradient of the loss \( \mathcal{L} \) with respect to \( W_1 \) is:
\[
\frac{\partial \mathcal{L}}{\partial W_1} = \frac{\partial \mathcal{L}}{\partial \mathbf{h}} \odot \sigma'(\mathbf{z}) \cdot \mathbf{x}^T.
\]

Here:
\begin{itemize}
    \item \( \frac{\partial \mathcal{L}}{\partial \mathbf{h}} \) is the gradient of the loss with respect to the hidden layer activations \( \mathbf{h} \),
    \item \( \sigma'(\mathbf{z}) \) is the derivative of the activation function \( \sigma \) applied element-wise to \( \mathbf{z} \),
    \item \( \mathbf{x}^T \) is the transpose of the input vector \( \mathbf{x} \).
\end{itemize}

\subsubsection{Where Does the Hadamard Product Come From?}
The key is the chain rule. The hidden layer activations \( \mathbf{h} \) are computed as:
\[
\mathbf{h} = \sigma(\mathbf{z}),
\]
where \( \mathbf{z} = W_1 \mathbf{x} + \mathbf{b}_1 \).

To compute \( \frac{\partial \mathcal{L}}{\partial W_1} \), we need to apply the chain rule through the activation function \( \sigma \). Specifically:
\[
\frac{\partial \mathcal{L}}{\partial W_1} = \frac{\partial \mathcal{L}}{\partial \mathbf{h}} \cdot \frac{\partial \mathbf{h}}{\partial \mathbf{z}} \cdot \frac{\partial \mathbf{z}}{\partial W_1}.
\]

\begin{itemize}
    \item \( \frac{\partial \mathbf{h}}{\partial \mathbf{z}} \) is the derivative of the activation function \( \sigma \) applied element-wise. Since \( \sigma \) is applied element-wise, \( \frac{\partial \mathbf{h}}{\partial \mathbf{z}} \) is a diagonal matrix with entries \( \sigma'(z_i) \). In practice, this is represented as a vector \( \sigma'(\mathbf{z}) \), and the multiplication becomes a Hadamard product (element-wise multiplication).

    \item \( \frac{\partial \mathbf{z}}{\partial W_1} \) is the derivative of \( \mathbf{z} \) with respect to \( W_1 \), which is simply \( \mathbf{x}^T \) because \( \mathbf{z} = W_1 \mathbf{x} + \mathbf{b}_1 \).
\end{itemize}

Thus, the chain rule gives:
\[
\frac{\partial \mathcal{L}}{\partial W_1} = \frac{\partial \mathcal{L}}{\partial \mathbf{h}} \odot \sigma'(\mathbf{z}) \cdot \mathbf{x}^T.
\]

\subsubsection{Why is \( \frac{\partial \mathbf{h}}{\partial \mathbf{z}} \) a Diagonal Matrix?}
The derivative \( \frac{\partial \mathbf{h}}{\partial \mathbf{z}} \) is a diagonal matrix because the activation function \( \sigma \) is applied \textbf{element-wise} to the vector \( \mathbf{z} \). This means that each element of the output vector \( \mathbf{h} \) depends only on the corresponding element of the input vector \( \mathbf{z} \). Mathematically, this can be written as:
\[
\mathbf{h} = \sigma(\mathbf{z}) = \begin{bmatrix} \sigma(z_1) \\ \sigma(z_2) \\ \vdots \\ \sigma(z_m) \end{bmatrix},
\]
where \( \mathbf{z} = \begin{bmatrix} z_1 \\ z_2 \\ \vdots \\ z_m \end{bmatrix} \).

The derivative \( \frac{\partial \mathbf{h}}{\partial \mathbf{z}} \) is a matrix that describes how each element of \( \mathbf{h} \) changes with respect to each element of \( \mathbf{z} \). Since \( \sigma \) is applied element-wise, the derivative of \( h_i \) with respect to \( z_j \) is zero when \( i \neq j \). This is because \( h_i \) depends only on \( z_i \) and not on any other \( z_j \) for \( j \neq i \).

Thus, the derivative matrix \( \frac{\partial \mathbf{h}}{\partial \mathbf{z}} \) will have non-zero entries only on the diagonal, where \( i = j \). The diagonal entries are the derivatives of \( \sigma(z_i) \) with respect to \( z_i \), which are \( \sigma'(z_i) \). Therefore, we can write:
\[
\frac{\partial \mathbf{h}}{\partial \mathbf{z}} = \begin{bmatrix}
\sigma'(z_1) & 0 & \cdots & 0 \\
0 & \sigma'(z_2) & \cdots & 0 \\
\vdots & \vdots & \ddots & \vdots \\
0 & 0 & \cdots & \sigma'(z_m)
\end{bmatrix}.
\]

This diagonal matrix is why the Hadamard product (element-wise multiplication) appears in the gradient computation for \( W_1 \).

\subsubsection{Explicit Development of the Hadamard Product}
Let’s make this explicit with an example. Suppose:
\begin{itemize}
    \item \( \frac{\partial \mathcal{L}}{\partial \mathbf{h}} = \begin{bmatrix} \frac{\partial \mathcal{L}}{\partial h_1} \\ \frac{\partial \mathcal{L}}{\partial h_2} \end{bmatrix} \),
    \item \( \sigma'(\mathbf{z}) = \begin{bmatrix} \sigma'(z_1) \\ \sigma'(z_2) \end{bmatrix} \),
    \item \( \mathbf{x}^T = \begin{bmatrix} x_1 & x_2 & x_3 \end{bmatrix} \).
\end{itemize}

The Hadamard product \( \frac{\partial \mathcal{L}}{\partial \mathbf{h}} \odot \sigma'(\mathbf{z}) \) is:
\[
\frac{\partial \mathcal{L}}{\partial \mathbf{h}} \odot \sigma'(\mathbf{z}) = \begin{bmatrix} \frac{\partial \mathcal{L}}{\partial h_1} \sigma'(z_1) \\ \frac{\partial \mathcal{L}}{\partial h_2} \sigma'(z_2) \end{bmatrix}.
\]

This vector is then multiplied by \( \mathbf{x}^T \) to compute the gradient \( \frac{\partial \mathcal{L}}{\partial W_1} \):
\[
\frac{\partial \mathcal{L}}{\partial W_1} = \begin{bmatrix} \frac{\partial \mathcal{L}}{\partial h_1} \sigma'(z_1) \\ \frac{\partial \mathcal{L}}{\partial h_2} \sigma'(z_2) \end{bmatrix} \cdot \begin{bmatrix} x_1 & x_2 & x_3 \end{bmatrix}.
\]

The result is a \( 2 \times 3 \) matrix, which matches the dimensions of \( W_1 \):
\[
\frac{\partial \mathcal{L}}{\partial W_1} = \begin{bmatrix}
\frac{\partial \mathcal{L}}{\partial h_1} \sigma'(z_1) x_1 & \frac{\partial \mathcal{L}}{\partial h_1} \sigma'(z_1) x_2 & \frac{\partial \mathcal{L}}{\partial h_1} \sigma'(z_1) x_3 \\
\frac{\partial \mathcal{L}}{\partial h_2} \sigma'(z_2) x_1 & \frac{\partial \mathcal{L}}{\partial h_2} \sigma'(z_2) x_2 & \frac{\partial \mathcal{L}}{\partial h_2} \sigma'(z_2) x_3
\end{bmatrix}.
\]

\subsubsection{Key Takeaways:}
\begin{itemize}
    \item The \textbf{Hadamard product} arises because the activation function \( \sigma \) is applied element-wise to the pre-activation values \( \mathbf{z} \).
    \item The derivative of \( \mathbf{h} \) with respect to \( \mathbf{z} \) is a diagonal matrix (or a vector of derivatives \( \sigma'(z_i) \)), which leads to element-wise multiplication.
    \item The chain rule combines this with the gradient of the loss \( \mathcal{L} \) with respect to \( \mathbf{h} \), resulting in the Hadamard product.
\end{itemize}

This explanation should clarify the origin of the Hadamard product in the gradient computation for \( W_1 \).

\subsection{Jacobian Matrix in Backpropagation}
The Jacobian matrix plays a key role in backpropagation when dealing with vector-valued functions. For example, the gradient of the loss \( \mathcal{L} \) with respect to the hidden layer \( \mathbf{h} \) is:
\[
\frac{\partial \mathcal{L}}{\partial \mathbf{h}} = W_2^T \cdot \frac{\partial \mathcal{L}}{\partial \mathbf{y}}.
\]
Here, \( W_2^T \) is the transpose of the weight matrix \( W_2 \), and it acts as the Jacobian matrix for the linear transformation \( \mathbf{y} = W_2 \mathbf{h} + \mathbf{b}_2 \).

\subsection{Recursive Application of the Chain Rule}
For a network with \( L \) layers, the gradients are computed recursively:
\begin{itemize}
    \item Gradients for the output layer:
    \[
    \frac{\partial \mathcal{L}}{\partial W_L} = \frac{\partial \mathcal{L}}{\partial \mathbf{y}_L} \cdot \mathbf{h}_{L-1}^T
    \]
    \[
    \frac{\partial \mathcal{L}}{\partial \mathbf{b}_L} = \frac{\partial \mathcal{L}}{\partial \mathbf{y}_L}
    \]
    \item For each hidden layer \( l = L-1, L-2, \dots, 1 \):
    \[
    \frac{\partial \mathcal{L}}{\partial \mathbf{h}_l} = W_{l+1}^T \cdot \frac{\partial \mathcal{L}}{\partial \mathbf{y}_{l+1}}
    \]
    \[
    \frac{\partial \mathcal{L}}{\partial W_l} = \frac{\partial \mathcal{L}}{\partial \mathbf{h}_l} \odot \sigma'(\mathbf{z}_l) \cdot \mathbf{h}_{l-1}^T
    \]
    \[
    \frac{\partial \mathcal{L}}{\partial \mathbf{b}_l} = \frac{\partial \mathcal{L}}{\partial \mathbf{h}_l} \odot \sigma'(\mathbf{z}_l)
    \]
\end{itemize}

\section{Symbol-to-Symbol Derivatives}
Symbol-to-symbol differentiation refers to computing derivatives using computational graphs, where each node represents a mathematical operation. Frameworks like TensorFlow and PyTorch use this approach.

For example, consider the function \( f(x, y) = x \cdot y + \sin(x) \). Its computational graph looks like:
\begin{itemize}
    \item Multiply \( x \) and \( y \).
    \item Compute \( \sin(x) \).
    \item Add the results of steps 1 and 2.
\end{itemize}

The derivatives are computed by traversing this graph backward, applying the chain rule at each step.

\section{General Back-Propagation}
General backpropagation extends the idea to any directed acyclic graph (DAG) representing a computation. The steps are:
\begin{itemize}
    \item Perform a forward pass to compute all intermediate values.
    \item Perform a backward pass to compute gradients using the chain rule.
\end{itemize}

This approach is used in modern deep learning frameworks to handle complex architectures like CNNs, RNNs, and transformers.

\section{Summary}
\begin{itemize}
    \item The \textbf{chain rule} is the backbone of backpropagation.
    \item The \textbf{Jacobian matrix} is used to compute gradients for vector-valued functions.
    \item The \textbf{gradient} of a scalar-valued function can be defined as a row vector or a column vector, depending on the convention. In this document, we use the row vector convention.
    \item \textbf{Recursive application} of the chain rule allows us to compute gradients efficiently.
    \item In \textbf{fully connected MLPs}, gradients are computed layer-by-layer.
    \item \textbf{Symbol-to-symbol differentiation} uses computational graphs to automate gradient computation.
    \item \textbf{General backpropagation} works for any DAG, enabling gradient computation in complex models.
\end{itemize}
\end{notechapter}
